\chapter{Conclusions and Future Work}
\label{ch:conclusion}

\section{Summary of Contributions}

This project conducted a comprehensive investigation of the Kaggle Hedge Fund Time Series Forecasting competition, progressing through three major phases:

\subsection{Exploratory Data Analysis}

The EDA provided the foundation for all subsequent modeling decisions:
\begin{itemize}
    \item Characterized the panel structure: 36,923 series across 23 entities, 4 horizons, with heterogeneous lengths.
    \item Identified the heavy-tailed target distribution with near-zero median.
    \item Discovered extreme weight concentration: top 1\% of rows $\to$ 65\% of the competition score.
    \item Established that \textbf{100\% of eligible series achieve stationarity after first differencing}.
    \item Found that individual feature correlations with the target are uniformly weak ($|\rho| < 0.10$).
    \item Determined that panel-wide seasonality is absent.
\end{itemize}

\subsection{Classical Modeling}

The classical modeling phase:
\begin{itemize}
    \item Evaluated 41 model specifications across 6 families (Baselines, AR, MA, ARMA, ARIMA, SARIMA).
    \item Empirically validated the EDA's stationarity finding through $d=0$ vs. $d=1$ head-to-head comparisons.
    \item Demonstrated that simple baselines (Expanding Mean, Rolling Mean) are surprisingly competitive.
    \item Produced per-series winner tables showing substantial heterogeneity across the panel.
\end{itemize}

\subsection{Deep Sequence Modeling}

The deep learning phase:
\begin{itemize}
    \item Implemented three recurrent architectures: vanilla RNN, GRU, and LSTM.
    \item Employed recursive forecasting for multi-step prediction.
    \item Developed adaptive training splits for series with limited pre-cutoff history.
    \item Demonstrated that gated architectures (GRU, LSTM) generally outperform vanilla RNN.
\end{itemize}

\section{Key Takeaways}

\begin{enumerate}
    \item \textbf{Financial time series forecasting is inherently difficult.} The efficient-market hypothesis implies that easily exploitable patterns are quickly arbitraged away. Our results confirm that even well-specified models struggle to substantially outperform simple baselines.

    \item \textbf{EDA-driven modeling is essential.} The stationarity analysis directly shaped the classical model specifications, and the weight concentration analysis informed metric selection. Modeling without thorough EDA would have produced misleading results.

    \item \textbf{Model diversity enables ensemble opportunities.} The per-series winner analysis reveals that no single model is universally best. An ensemble that selects or blends models per-series could capture the best of each approach.

    \item \textbf{The weight column is a first-class design element.} Its extreme concentration means that competition performance is dominated by a small fraction of observations. Weight-aware training and evaluation are non-negotiable.
\end{enumerate}

\section{Limitations}

\begin{enumerate}
    \item \textbf{Univariate deep models:} The current deep learning implementation uses only the target variable, not the 86 predictor features. This leaves substantial potential signal unexploited.

    \item \textbf{Per-series training:} Both classical and deep models are trained independently per series, ignoring cross-series structure that panel methods could exploit.

    \item \textbf{No weighted training loss:} The deep models use unweighted MSE loss, which does not align with the competition's weighted evaluation metric.

    \item \textbf{Limited hyperparameter tuning:} Due to computational constraints, the deep learning hyperparameters were not exhaustively optimized.

    \item \textbf{No gradient-boosted tree models:} LightGBM, XGBoost, and CatBoost---commonly dominant in tabular forecasting competitions---were not implemented in this study.
\end{enumerate}

\section{Future Work}

Several promising directions remain:

\subsection{Feature-Based Models}
Incorporate the 86 predictor features using:
\begin{itemize}
    \item Gradient-boosted trees (LightGBM, XGBoost) with engineered lag features.
    \item Transformer-based architectures (e.g., Temporal Fusion Transformer) that can attend to both temporal and cross-sectional patterns.
    \item Multi-input recurrent networks that process feature vectors alongside the target sequence.
\end{itemize}

\subsection{Weighted Loss Functions}
Implement competition-aligned training by:
\begin{itemize}
    \item Using weighted MSE during deep model training.
    \item Developing differentiable approximations to the weighted skill score for end-to-end optimization.
\end{itemize}

\subsection{Ensemble Methods}
Combine models using:
\begin{itemize}
    \item Per-series model selection based on validation performance.
    \item Stacking with a meta-learner trained on validation predictions.
    \item Bayesian model averaging weighted by per-series posterior probabilities.
\end{itemize}

\subsection{Global and Panel Models}
Move beyond per-series training:
\begin{itemize}
    \item Train a single model across all series with entity embeddings (similar to N-BEATS or DeepAR).
    \item Use hierarchical models that share parameters across related series while allowing series-specific adjustments.
\end{itemize}

\subsection{Advanced Architectures}
Explore state-of-the-art time series models:
\begin{itemize}
    \item \textbf{Transformer models:} Self-attention can capture long-range dependencies without recurrence.
    \item \textbf{N-BEATS:} Interpretable deep learning architecture designed for time series forecasting.
    \item \textbf{Neural basis expansion:} Models that decompose forecasts into interpretable trend and seasonality components.
\end{itemize}

\section{Closing Remarks}

This project demonstrates that rigorous financial time series forecasting requires a multi-faceted approach: thorough exploratory analysis, principled baseline establishment, and progressive model sophistication. While no single approach achieves dominant performance across the heterogeneous panel, the systematic methodology provides a solid foundation for further improvement through ensemble methods and feature-based modeling.

The competition's design---with anonymized data, weighted evaluation, and strict temporal splitting---mirrors the challenges of real-world quantitative finance, where model robustness and out-of-sample generalization are paramount.
